% Unbounded local boundary comparison and its global layer consequence.
\section{Length-sensitive reflection without a long run of ones}
For an increasing arm $Y$ of weight $b$ and length $L$, put
\[
 F_Y=\sum f_jz^j,\quad q_Y=z^bF_Y(1/z),\quad
 C_Y(r)=\sum_{j=0}^r(q_Y)_j,
\]
with $C_Y(r)=0$ for $r<0$. Use the formal power series
\[
 B_Y=1+\left(\frac{z}{1-z}-1\right)F_Y.
\]
It counts the minimal failures obtained by prefixing a positive
letter to an embedded word. Its first nonzero coefficient is at
the minimum missing weight $\kappa(Y)$; its coefficients are
nonnegative. For any finite peak $M$, its coefficients through
rank $M$ agree with those of $D_M(Y)$.

\begin{theorem}[A local reverse-boundary comparison]
For every increasing arm, including the empty arm with $\kappa=1$,
\[
 C_Y(r)\ge (B_Y)_{\kappa+r-1}+1
 \qquad(0\le r\le L-\kappa+1).
\]
Consequently, if $0\le h\le\kappa-1$, $s\ge0$, and $h+s\le L$,
then
\[
 C_Y(s)\ge(B_Y)_{h+s}+1.
\]
\end{theorem}
\begin{proof}
For the empty arm the sole comparison is $1\ge0+1$.
For $Y=1^L$, $L\ge1$, one has $\kappa=2$,
$C_Y(r)=r+1$ and $(B_Y)_{r+1}=r$ throughout $0\le r\le L-1$.
These give the initial cases, after which only parts $t\ge2$
need be appended. For a prefix of length $L$, weight $c$, and
missing weight $k$, the exact recurrences are
\[
 C_{Yt}(r)=\sum_{j=0}^{t-1}C_Y(r-j)+\mathbf1_{\{r\ge c+t\}},
 \qquad B_{Yt}=S_tB_Y+\frac{z^{t+1}}{1-z}.
\]
Also $k'=\min(k+1,t+1)$ and $k\le t+1$.

If $k'=k+1$, the new admissible range is $r\le L-k+1$,
exactly the old range. The difference in the desired comparison is
\[
 \sum_{j=0}^{t-1}
   \bigl(C_Y(r-j)-(B_Y)_{k+r-j-1}\bigr)
   -\mathbf1_{\{k+r\ge t+1\}}.
\]
The extra indicator from the cumulative recurrence vanishes here:
$r\le L$ and $c+t\ge L+2$. Terms with negative $r-j$ vanish,
and every other term is at least one. If the subtracted unit
occurs, $t\ge k$ forces $r\ge1$, so $t\ge2$ supplies at
least two such terms. The result is at least one.

If $k'=k=t+1$, the new range is $r\le L-k+2$, one index
longer than the old one. The summands now compare
$C_Y(s)$ with $(B_Y)_{k+s-2}$, where $s=r-j$.
For $s=0$ their difference is one. For $s\ge1$, use the old
comparison at $s-1\le L-k+1$ and cumulative monotonicity:
\[
 C_Y(s)\ge C_Y(s-1)\ge(B_Y)_{k+s-2}+1.
\]
Negative $s$ again gives zero. The subtracted unit occurs only
when $r\ge1$, and again at least two strict summands pay for it.
This proves the first assertion by induction.

For the consequence, if $h+s<\kappa$ the boundary coefficient
is zero and $C_Y(s)\ge1$. Otherwise put
$r=h+s-\kappa+1$. Then $0\le r\le s$ and
$r\le L-\kappa+1$. Apply the first assertion and monotonicity
of the cumulative sequence.
\end{proof}

\begin{proposition}[The next boundary rank and the low height layers]
The non-strict inequality
\[
 C_Y(r)\ge(B_Y)_{\kappa+r-1}
\]
holds on the one-index longer range $0\le r\le L-\kappa+2$.
Consequently $C_Y(s)\ge(B_Y)_{h+s}$ whenever
$0\le h\le\kappa-1$, $s\ge0$, and $h+s\le L+1$.
For a strict peak $M$, every height layer of $D_M(Y)$ whose
left endpoint is at most
$\min\{L+1,M+\kappa-2\}$ has endpoint sum at least
$b+M+\kappa-1$.
\end{proposition}
\begin{proof}
Only the additional index needs proof. For an all-one arm it
is equality. Induct through non-one appends as in the local
theorem. When the missing weight increases, the first summand
at the new last index uses the old non-strict boundary inequality;
all the remaining nonnegative-index summands use the strict
inequality. The new last index is at least one, so the second
summand exists and is strict. It pays the possible subtracted unit.

When the missing weight stays $k=t+1$, shift the old comparison
index down by one, exactly as in the strict proof. The first
summand then uses the old non-strict boundary comparison, and
the second uses the old strict comparison. The new last index
is at least two, so the second term again exists and pays the
unit. The extra cumulative recurrence term still vanishes,
since the comparison index is at most $L+1<c+t$.
This proves the extended inequality. The assertion with $h$
follows by cumulative monotonicity, as before.

For a layer with left endpoint $l$ in the stated range, put
$s=l-\kappa+1$. Then $0\le s<M$, so the reverse kernel
coefficient at $s$ is $C_Y(s)$. Also $l\le L+1$ and
$(D_M(Y))_l\le(B_Y)_l$, with equality when $l\le M$.
The extended comparison therefore gives
$\delta_s\ge(D_M(Y))_l$. The layer reaches at least
$b+M-s$, and its endpoint sum is at least $b+M+\kappa-1$.
\end{proof}

\begin{lemma}[One more layer offset after a sufficiently long prefix]
Let the increasing strict prefix $Y$ have length $L\ge M-2$,
weight $c$, and missing weight $k$. Suppose every layer of
$D_M(Y)$ has endpoint sum at least $c+M+h$, where
$1\le h\le k-1$. Appending an increasing part $2\le t<M$
raises this bound to $c+t+M+h+1$.
\end{lemma}
\begin{proof}
Write $b=c+t$, $\delta=z^{c+M}D_M(Y)(1/z)$, and $W=[t]\delta$.
The old layer bound makes every layer of $W$ have endpoint sum
at most $C=b+M-h-1$. The new reverse kernel is
$W+z^b[M-t]$. The unit interval is strictly to the right of
$C/2$, since $b\ge M$.

For a corrected index $j\in[b,b+M-t-1]$, its mirror is
$i=C-j\in[t-h,M-h-1]$. It is nonnegative because
$h\le k-1\le t$. Direct reversal of the convolution gives
\[
 W_i-W_j=\sum_{v=0}^{t-1}
    \bigl(\delta_{i-v}-(D_M(Y))_{h+i-v}\bigr).
\]
For $s=i-v\ge0$, one has $s<M$ and $h+s\le M-1\le L+1$.
Thus $\delta_s=C_Y(s)$, and the local theorem makes the
corresponding summand nonnegative, and at least one if $h+s\le L$.
For negative $s$ both coefficients vanish, since $h+s<k$.
If the $v=0$ term is not covered by the strict range, then
$L=M-2$ and $h+i=M-1=L+1$. Since $h\le k-1\le L$,
one has $i\ge1$; the $v=1$ term exists because $t\ge2$,
and it is strict. In every case at least one term is strict.
This pays for the corrected unit. All other reflected
comparisons are inherited from $W$. The new kernel is already
known to be unimodal, so these comparisons give its layer bound.
\end{proof}

\begin{corollary}[A richer short prefix improves the bound by one]
Let $M\ge4$ and suppose the prefix of $Y$ consisting of its first
$M-2$ parts has weight at least $M$. Then, for $L\ge M-2$,
every layer of $D_M(Y)$ has endpoint sum at least
\[
 b+M+\min\{\kappa(Y)-1,L-M+4\}.
\]
\end{corollary}
\begin{proof}
The prefix has the known offset two, and its missing weight is
at least three. Starting at length $M-2$, apply the strengthened
append lemma and the same minimum-update induction.
\end{proof}

\begin{theorem}[A length-sensitive global layer bound]
Let $M\ge3$, let $Y$ be an increasing strict arm of length
$L\ge M-1$, weight $b$, and missing weight $\kappa$. Then
every layer of $D_M(Y)$ has endpoint sum at least
\[
 b+M+\min\{\kappa-1,\ L-M+4\}.
\]
Within the stated length range, the full offset $\kappa-1$ holds
whenever $L\ge M+\kappa-5$. The simpler sufficient condition
$L\ge\max(M-1,2M-5)$ works for every arm.
\end{theorem}
\begin{proof}
For $M=3$, a nonrectangular arm of length at least two has weight
at least three and $\kappa=3$, so the known offset two suffices;
all-one arms use the earlier long-one theorem. Assume $M\ge4$
and inspect the first $M-2$ parts. If their weight is at least
$M$, the richer-prefix corollary proves the result.

If their weight is $M-2$, they are all ones. Either the complete
arm starts with at least $M-1$ ones, or its first non-one part
$t$ completes a two-factor seed $1^{M-2}t$ of weight at least
$M$. The earlier seed theorems give the full offset in either case.

The remaining prefix has weight $M-1$ and is exactly
$1^{M-3}2$. Let $t$ be the next part. If $t=2$, the resulting
prefix has weight $M+1$, missing weight three, and hence the
known full offset two. If $t\ge3$, it has missing weight four
and the full offset three, by the explicit seed calculation below.
Once a prefix of length at least $M-1$ has the full offset
$\kappa-1$, the append lemma preserves that assertion:
it raises the offset to $\kappa$, and the new missing weight
is at most $\kappa+1$. This handles every continuation.
Finally $\kappa\le M$ gives the uniform length condition.
\end{proof}

\begin{lemma}[The remaining short seed]
For $M\ge4$ and $3\le t<M$, the kernel of
$Y=1^{M-3}2t$ has every layer endpoint sum at least $b+M+3$.
\end{lemma}
\begin{proof}
Here $b=M-1+t$, $\kappa=4$, and put $C=2M+t-4$.
The reverse kernel is
\[
 \delta=[2][t][M-1][M-2]
       +z^{M-1}[t]([2]+[M-2])
       +z^{M-1+t}[M-t].
\]
The first box product has degree $C-1$. After subtracting its
reflection about $C/2$ and simplifying the other summands, one gets
\[
 \delta-z^C\delta(1/z)
 =(1-z^M)[t][M-2]
 -z^{M-3}(1-z^2)[2][t]
 +(z^{M+t-1}-z^{t-2})[M-t].
\]
Let $c_i=[z^i][t][M-2]$. At every left index $2i\le C$, the
coefficient of this difference is exactly
\[
 c_i-c_{i-M}
 -\mathbf1_{\{t-2\le i\le M-3\}}
 -\mathbf1_{\{i=M-3\}}
 -2\mathbf1_{\{i=M-2\}}
 -\mathbf1_{\{i=M-1\}}.
\]
For $i\le M-4$, any subtracted unit is paid by $c_i\ge1$.
At $i=M-3,M-2,M-1$, the box coefficients are respectively
at least $2,2,1$, paying all displayed costs. At $i=M+j$
on the left half, the difference is $t-4-2j\ge0$ because
$2j\le t-4$. This proves reflection and hence the asserted
layer bound, using the established kernel unimodality.
\end{proof}

\paragraph{An insufficient shorter threshold.}
Length $M-1$ by itself does not force the full offset
$\kappa-1$. For $M=9$ and $Y=(1,1,2,3,4,5,8,8)$, one has
$L=8$, $\kappa=8$, and the minimum layer endpoint sum minus
$b+M$ is six. The asserted minimum in the theorem is essential;
the local strict comparison alone does not imply global reflection.

\begin{theorem}[A padding criterion using length and missing weight]
Let $X=1^\alpha U$, $Y=1^\beta V$ be increasing strict arms,
$u=\operatorname{wt}U$, $s=\operatorname{len}V$,
$b=\operatorname{wt}Y$, and $L=\beta+s$.
For $M\ge3$ define
\[
 h_M(Y)=
 \begin{cases}
  \min\{\kappa(Y)-1,L-M+4\},&L\ge M-1,\\
  2,&L=M-2\text{ and }b\ge M,\\
  1,&\text{otherwise},
 \end{cases}
\]
If either orientation satisfies
\[
 \alpha\ge u+s+\max(\beta-M,0)-h_M(Y)-1,
\]
then both relative polynomials and the full core polynomial
are unimodal. Arms of non-one weight at most four remain
unconditional base cases by the earlier theorem.
\end{theorem}
\begin{proof}
The layer bound gives relative growth through
$g=\lfloor(a+b+M+h_M(Y)+1)/2\rfloor$.
The grouped-one mode bound gives
$m_Y\ge\lfloor(b-s+M-1-r)/2\rfloor$, where
$r=\max(\beta-M,0)$. The displayed condition implies
$g+m_Y\ge u+b+M-1$, so the terminal coefficient bridge applies.

For the least-counterexample argument, the qualifying family is
closed under final-part deletion and arm exchange after adjoining
the proved small-skeleton base cases. Deleting a
non-one part of $U$ decreases $u$. Deleting one from $V$
decreases $s$ by one and decreases $h_M(Y)$ by at most one,
unless the shorter arm is already an unconditional base case:
within the longer-length branch this follows because both
arguments of the minimum decrease by at most one. At the crossing
to $L=M-2$, the old offset is at most three. If the shorter arm
has weight at least $M$, its offset is two. Otherwise an increasing
arm of length $M-2$ and weight below $M$ is all ones or
$1^{M-3}2$, both unconditional base cases. The crossing from
$L=M-2$ to $L=M-3$ loses at most one. Therefore the right-hand side
of the padding condition cannot increase. The usual balanced-band
equivalence gives both relative orientations and the full core.
This is a sufficient family, not an unrestricted solution.
\end{proof}
